\documentclass[10pt,a4paper]{amsart}
\usepackage[margin=19mm]{geometry}
\usepackage{amsmath,amssymb,mathtools}
\usepackage[scr=rsfs,cal=euler]{mathalfa}
\usepackage{xfrac}
\usepackage[unicode,hidelinks,bookmarks=false]{hyperref}
\newcommand{\ie}{i.e.\ }
\newcommand{\cf}{cf.\ }
\newcommand{\resp}{respectively\ }
\newcommand{\Subsection}{\subsection}
\providecommand{\nobreakdash}{\nobreak\hbox{-}}
\setlength{\emergencystretch}{2em}
\multlinegap0pt
\usepackage{bm}
\usepackage{bbm}
\usepackage{dsfont}
\usepackage{xspace}
\usepackage{cancel}
\usepackage{mathtools}
\usepackage{amsthm}
\makeatletter
\@ifundefined{Th}{\newtheorem{Th}{Th}}{}
\@ifundefined{Lemma}{\newtheorem{Lemma}[Th]{Lemma}}{}
\@ifundefined{Prop}{\newtheorem{Prop}[Th]{Proposition}}{}
\makeatother
\newcommand{\haupt}{\lambda}
\newcommand{\upee}[1]{\left. #1 \right\rvert \hspace{1pt} U_{p}}
\newcommand{\utwo}[1]{\left. #1 \right\rvert \hspace{1pt} U_{2}}
\title[An Improved Lower Bound on the Image of the 2-adic Character]{An Improved Lower Bound on the Image of the 2-adic Character Map for the Heisenberg Algebra via Modular Linear Differential Equations}
\author{Daniel Barake, Cameron Franc}
\date{Source: arXiv:2503.11817v2 (CC BY 4.0)}
\begin{document}
\maketitle

\noindent\emph{Source and licence.} An Improved Lower Bound on the Image of the 2-adic Character Map for the Heisenberg Algebra via Modular Linear Differential Equations. Version arXiv:2503.11817v2 (CC BY 4.0), distributed under the Creative Commons Attribution 4.0 International licence, as recorded in the arXiv licence metadata for this exact version and shown on the abstract page https://arxiv.org/abs/2503.11817v2. Full rights notice: licenses/arxiv-2503.11817-CC-BY-4.0.txt.\par\smallskip

\noindent\textit{Editorial scope.} Counted unit: the Prop whose source label is \texttt{Prop: hmdiffs} in the article (source0.tex lines 372--378, source label \texttt{Prop: hmdiffs}), delivered together with the article's own complete original proof of it (source0.tex lines 379--440, one \texttt{proof} environment of 62 lines). No mathematics is written, repaired or re-derived: statement and proof are the article's own text line for line. Required context retained uncounted: Lemma: Upoperator (lines 276--283); Lemma: hdiffeq (lines 316--321); hmsbetter (lines 365--367). Every definition, display and label of those lines is retained unchanged. Omitted from this package and not redistributed: 1--275, 284--315, 322--364, 368--371, 441--737. Nothing inside a retained excerpt is omitted or altered: every retained body is delivered exactly as the article's lines are written, so no omission receipt of any kind accompanies this package. Citations: the retained excerpts carry no citation command at all, so no bibliography entry of the article is transcribed and no reference is redistributed. Reference adaptation: none was needed and none was made; every reference inside the delivered unit resolves inside it, so no unresolved reference or citation is produced. Printed slips kept literal: no printed word, symbol, hypothesis or number of the counted statement or of its proof was altered. Layout and class: the article's own class is \texttt{amsart} with packages including amsthm, amsmath, mathtools, amscd, braket, fontenc, eucal, MnSymbol, parskip, geometry, tikz-cd; the wrapper is the lane's pinned \texttt{amsart} editorial wrapper at 10pt on A4, which restates the theorem-like environments the retained text needs and adds the article's own macro definitions that the retained text needs (\texttt{\textbackslash haupt}, \texttt{\textbackslash upee}, \texttt{\textbackslash utwo}), with extra wrapper packages (bm, bbm, dsfont, xspace, cancel, mathtools). The delivered numbering is the unit's own. Assets: no retained span contains a figure environment, a graphics-inclusion command, a PDF-page inclusion, a table or a TikZ picture; no image file is copied into this package, the package declares no asset dependency and no image file is redistributed with it. Licence: version arXiv:2503.11817v2 of this article is distributed under the Creative Commons Attribution 4.0 International licence, as recorded in the arXiv licence metadata for this exact version and shown on the abstract page https://arxiv.org/abs/2503.11817v2. A full rights notice is delivered at licenses/arxiv-2503.11817-CC-BY-4.0.txt. Characters: every retained excerpt is pure ASCII, so the delivered engine, XeLaTeX with its default Latin Modern text font, reports no missing character.
\begin{Lemma} \label{Lemma: Upoperator}
    Let $p$ be prime. For formal Laurent series $f(q), g(q) \in \mathbb{C}[[q]]$, we have the identities
    \begin{enumerate}
        \item $\upee{ \left( f(q) g(q^{p}) \right) } =  g (q) \upee{ \left( f (q) \right) }$
        \item $\upee{ \theta (f(q)) } = p \cdot \theta \left( \upee{ f(q) }  \right)$ where $\theta = q \frac{d}{dq}$
        \item For all $k \geq 2$ even, $\upee{ E_{k}^{\star}(q) }= E_{k}^{\star}(q)$.
    \end{enumerate}
\end{Lemma}

\begin{Lemma} \label{Lemma: hdiffeq}
Let $p=2$. For $n \geq 1$, $\Delta (q)^{n} \haupt^{n}$ satisfies the following third-degree MLDE:
\begin{align}
    \mathcal{D}^{3} \left( \Delta (q)^{n} \haupt^{n} \right) = \left( \left( n^{2} - n^{3} \right) E_{2}^{\star}(q)^{3} - \left( \frac{n^{2}}{2} + \frac{n}{18} \right) E_{2}^{\star}(q)E_{4}(q) \right) \Delta (q)^{n} \haupt^{n}.\label{hdiffeq}
\end{align}
\end{Lemma}

\begin{align}
    \haupt_{n,m} = \haupt^{n}  \left( 5E_{4}(q) - 20 E_{2}^{\star}(q)^{2} \right)^{3 \cdot 2^{m}} + 2^{1 - 12n} (240)^{3 \cdot 2^{m}} \Delta (q)^{2^{m+1}-n} \utwo{\left( \Delta (q)^{n-2^{m}} \right) }. \label{hmsbetter}
\end{align}

\begin{Prop} \label{Prop: hmdiffs}
    Set $p=2$. For $n,m \geq 1$, the modular form $\Delta (q)^{n} \haupt_{n,m}$ satisfies the third-degree MLDE:
    \begin{multline}
        \mathcal{D}^{3} ( \Delta (q)^{n} \haupt_{n,m}) = \left( \frac{3}{4}t_{n,m}^{2} + \frac{1}{4}t_{n,m} + \frac{1}{18} \right) \mathcal{D}(\Delta (q)^{n} \haupt_{n,m}) E_{4}(q) \\ - \frac{1}{4} \left( t_{n,m}^{3} + t_{n,m}^{2} \right) \Delta (q)^{n}\haupt_{n,m} E_{6}(q),
    \end{multline}
    where $t_{n,m} = 2^{m}-n$.
\end{Prop}

\begin{proof}
    As in the proof of Lemma \ref{Lemma: hdiffeq}, it suffices to show that the $\haupt_{n,m}$ satisfy the MLDE in the claim with the modular derivative acting with weight $4a \cdot 2^{m}$. Let us temporarily set
   \begin{align*}
       X = 5E_{4}(q) - 20 E_{2}^{\star}(q)^{2}
   \end{align*}
   so that eq. (\ref{hmsbetter}) becomes
   \begin{align*}
        \haupt_{n,m} &= \haupt^{n}  X^{3 \cdot 2^{m}} + 2^{1 - 12n} (240)^{3 \cdot 2^{m}} \Delta (q)^{2^{m+1}-n} \utwo{ \left( \Delta (q)^{n-2^{m}} \right) } \\
        &:= T_{1} + 2^{1-12n} (240)^{3 \cdot 2^{m}}  T_{2}.
   \end{align*}
   We show that $T_{1}$ and $T_{2}$ satisfy the MLDE in the statement of the proposition separately so that by linearity the result holds. Consider first $T_{1}$. It is easily verified that $\mathcal{D}(X) = (1/3)E_{2}^{\star}(q)X$, so that
   \begin{align}
       \mathcal{D}(T_{1}) = t_{n,m} \haupt^{n} E_{2}^{\star}(q) X^{3 \cdot 2^{m}} \label{T1d1}
   \end{align}
   and from this
   \begin{align*}
       \mathcal{D}^{2}(T_{1}) = \haupt^{n} X^{3 \cdot 2^{m}} \left( \left( t_{n,m}^{2} + \frac{1}{3} t_{n,m} \right) E_{2}^{\star}(q)^{2} - \frac{1}{6} t_{n,m} E_{4}(q) \right).
   \end{align*}
   Lastly, a lengthy computation shows that
   \begin{align*}
       \mathcal{D}^{3}(T_{1}) = \haupt^{n}X^{3 \cdot 2^{m}} \left( \left( t_{n,m}^{3} + t_{n,m}^{2} \right) E_{2}^{\star}(q)^{3} + \left( \frac{1}{18} t_{n,m} - \frac{1}{2} t_{n,m}^{2} \right) E_{2}^{\star}(q)E_{4}(q) \right).
   \end{align*}
   Now, using the identity $E_{6}(q) = -4 E_{2}^{\star}(q)^{3} + 3 E_{2}^{\star}(q) E_{4}(q)$, we can re-arrange
   \begin{multline*}
       \mathcal{D}^{3}(T_{1}) = \haupt^{n} X^{3 \cdot 2^{m}} \left( \left( \frac{3}{4} t_{n,m}^{3} + \frac{1}{4} t_{n,m}^{2} + \frac{1}{18} t_{n,m} \right) E_{2}^{\star}(q)E_{4}(q) \right. \\ \left. - \frac{1}{4}\left( t_{n,m}^{3} + t_{n,m}^{2} \right) \left( -4 E_{2}^{\star}(q)^{3} + 3 E_{2}^{\star}(q)E_{4}(q) \right) \right)
   \end{multline*}
   and so using (\ref{T1d1}) we get that $T_{1}$ satisfies the MLDE in the statement. Now for $T_{2}$: Since $\mathcal{D}_{1/2}(\eta(q)) = 0$, we have
   \begin{align*}
       \mathcal{D}^{i}(T_{2}) = \Delta (q)^{2^{m+1}-n} \mathcal{D}^{i} \left( \utwo{  \Delta (q)^{-t_{n,m}} } \right)
   \end{align*}
   for $i = 1,2,3$. In particular, using the fact that $\theta (\eta (q)) = (1/24) E_{2}(q) \eta (q)$ as well as Lemma \ref{Lemma: Upoperator} (2),
   \begin{align}
        \mathcal{D} \left. \left( \Delta (q)^{-t_{n,m}} \right\rvert \hspace{1pt} U_{2} \right) = - \frac{1}{2} t_{n,m} \utwo{ \left( \Delta (q)^{- t_{n,m}} E_{2}^{\star}(q) \right) }. \label{T2d1}
   \end{align}
   From here, we get
   \begin{align*}
       \mathcal{D}^{2} \left( \utwo{ \Delta (q)^{-t_{n,m}} } \right) = \utwo{ \left( \Delta (q)^{-t_{n,m}} \left( \left( \frac{ t_{n,m}^{2}}{4} - \frac{ t_{n,m}}{8} \right) E_{2}^{\star}(q)^{2} + \frac{ t_{n,m}}{24} E_{4}(q) \right) \right) },
   \end{align*}
   and 
   \begin{multline*}
       \mathcal{D}^{3} \left( \utwo{ \Delta (q)^{- t_{n,m}} } \right) = \left. \left( \Delta (q)^{- t_{n,m}} \left( \left( - \frac{t_{n,m}^{3}}{8} + \frac{3 t_{n,m}^{2}}{16} - \frac{5 t_{n,m}}{144} \right) E_{2}^{\star}(q)^{3} \right.\right. \right. \\
       + \utwo{ \left. \left. \left( - \frac{t_{n,m}^{2}}{16} + \frac{t_{n,m}}{144} \right) E_{2}^{\star}(q) E_{4}(q) \right) \right) }.
   \end{multline*}
   Making use of the identities
       \begin{align*}
           E_{4}(q^{2}) &= \frac{5}{4} E_{2}^{\star}(q)^{2} - \frac{1}{4} E_{4}(q) \\
           E_{6}(q^{2}) &= \frac{3}{8}E_{2}^{\star}(q)E_{4}(q) - \frac{11}{8}E_{2}^{\star}(q)^{3}
       \end{align*}
   alongside Lemma \ref{Lemma: Upoperator} (2), we can re-arrange
   \begin{multline*}
           \mathcal{D}^{3} \left( \utwo{ \Delta (q)^{-t_{n,m}} } \right) = \left. \left( \left( - \frac{15}{32} t_{n,m}^{3} - \frac{5}{32} t_{n,m}^{2} - \frac{5}{144} t_{n,m} \right) \Delta (q)^{-t_{n,m}}E_{2}^{\star}(q)^{3} \right. \right. \\ 
           \left. \left. + \left( \frac{3}{32} t_{n,m}^{3} + \frac{1}{32} t_{n,m}^{2} + \frac{1}{144} t_{n,m} \right) \Delta (q)^{- t_{n,m}} E_{2}^{\star}(q)E_{4}(q) \right. \right. \\
           \left. \left. - \frac{3}{32} \left( t_{n,m}^{3} + t_{n,m}^{2} \right) \Delta (q)^{- t_{n,m}} E_{2}^{\star}(q)E_{4}(q) \right. \right. \\
           \utwo{ \left. + \frac{11}{32} \left( t_{n,m}^{3} + t_{n,m}^{2} \right) \Delta (q)^{- t_{n,m}} E_{2}^{\star}(q)^{3} \right) }
       \end{multline*}
       so that by (\ref{T2d1})
       \begin{multline*}
           \mathcal{D}^{3} \left( \utwo{ \Delta (q)^{- t_{n,m}} } \right) = \left( \frac{3}{4} t_{n,m}^{2} + \frac{1}{4} t_{n,m} + \frac{1}{18} \right) \mathcal{D} \left( \utwo{ \Delta (q)^{- t_{n,m}} } \right) E_{4}(q) \\
           - \frac{1}{4} \left( t_{n,m}^{3} + t_{n,m}^{2} \right) \left( \utwo{ \Delta (q)^{- t_{n,m}} } \right) E_{6}(q)
       \end{multline*}
       which shows that $T_{2}$ satisfies the MLDE in the proposition.
\end{proof}

\end{document}
